\documentclass[11pt]{article}

\usepackage[margin=1in]{geometry}
\usepackage[T1]{fontenc}
\usepackage[utf8]{inputenc}
\usepackage{textcomp}
\usepackage{amsmath}
\usepackage{amssymb}
\usepackage{booktabs}
\usepackage{array}
\usepackage{calc}
\usepackage{longtable}
\usepackage{xcolor}
\usepackage{algorithm}
\usepackage[noend]{algpseudocode}
\usepackage[hidelinks]{hyperref}

\setlength{\emergencystretch}{3em}
\setlength{\parskip}{0.35em plus 0.1em minus 0.05em}
\setcounter{secnumdepth}{0}
\providecommand{\tightlist}{%
  \setlength{\itemsep}{0pt}\setlength{\parskip}{0pt}}
\ifdefined\DeclareUnicodeCharacter
  \DeclareUnicodeCharacter{00A7}{\S}
\fi
\newcommand{\R}{\mathbb{R}}

\title{Appendix B\\Model Specifications for an ABM of Brokerage in Matching Markets}
\author{}
\date{}

\begin{document}
\maketitle

This document specifies an agent-based model of brokered matching.

The model is a discrete-time simulation of a matching market with \emph{agents} and a \emph{broker}. Agents seeking pairwise matches either search on their own or outsource to the broker. All economic quantities (match output, fees, surplus) are in the same monetary units. All agents use heuristic decision rules. No agent holds beliefs about other agents' strategies.

The model specification covers agents (§1), the matching problem (§2), how agents learn to predict match quality (§3), match economics (§4), network structure and agent turnover (§5), search and broker access (§6), the outsourcing decision (§7), and period measures (§8). Baseline parameter values and reproducibility are documented in Part II (§9).

\section{Part I. Model Specification}\label{part-i.-model-specification}

\subsection{1. Agents and Time}\label{agents-and-time}

The model runs for \(T\) discrete periods (default 500). It has \(N\) agents (default 1000), a single broker, and an undirected network \(G\) over the agents and broker.

Each agent \(i\) is characterized by:

\begin{itemize}
\tightlist
\item \textbf{Type} \(\mathbf{x}_i \in \mathbb{R}^d\): a fixed vector of observable characteristics, drawn as specified below and used by the matching function (§2).
\item \textbf{Current-period matches} \(M_i^t\): active bilateral relationships involving \(i\) in period \(t\). Each counterparty can appear at most once in \(M_i^t\); incoming accepted offers are not capped.
\item \textbf{Active demand bound} \(K\): the maximum number of active demands an agent can initiate in a period (default \(K = 5\); §6).
\item \textbf{Experience history} \(\mathcal{H}_{i}^t = \{(\mathbf{x}_j, q_{ij})\}\): (other party's type, realized match output) pairs from all matches involving \(i\) (§3a).
\item \textbf{Satisfaction indices} \(s_{i,c}^t\): one scalar per search channel \(c \in \{\text{self}, \text{broker}\}\), used in the outsourcing decision (§7).
\end{itemize}

\subsubsection{Agent Types}\label{agent-types}

Type vectors are generated near a smooth one-dimensional curve on the unit sphere in \(\mathbb{R}^d\). The curve is parameterized by a position \(u \in [0, 1]\):

\[\mathbf{x}(u) = \frac{\tilde{\mathbf{x}}(u)}{\|\tilde{\mathbf{x}}(u)\|}, \qquad \tilde{x}_\ell(u) = \begin{cases} \sin(2\pi \nu_\ell u + \psi_\ell) & \ell = 1, \ldots, d_\gamma \\ 0 & \ell = d_\gamma+1, \ldots, d \end{cases}\]

where \(\nu_\ell \sim \mathrm{Unif}\{1, 2, 3, 4, 5\}\) are random integer frequencies, \(\psi_\ell \sim \mathrm{Unif}[0, 2\pi)\) are random phases, both drawn once per simulation, and \(d_\gamma \leq d\) is the number of \textbf{active dimensions}. The remaining dimensions enter only through noise.

Each agent is drawn at a random position \(u_i \sim \mathrm{Unif}[0,1]\) on the curve, then perturbed:

\[\mathbf{x}_i = \frac{\mathbf{x}(u_i) + \boldsymbol{\epsilon}_i}{\|\mathbf{x}(u_i) + \boldsymbol{\epsilon}_i\|}, \qquad \boldsymbol{\epsilon}_i \sim \mathcal{N}\!\left(\mathbf{0}, \frac{\sigma_x^2}{d} \mathbf{I}_d\right)\]

Noise is applied in all \(d\) dimensions, then the vector is re-normalized. The scale \(\sigma_x / \sqrt{d}\) keeps the expected displacement from the curve approximately \(\sigma_x\) for any \(d\).

\(d_\gamma\) controls type-space complexity by setting how many dimensions carry curve signal.

\subsubsection{Broker}\label{broker}

The broker is a permanent node in \(G\) and has no type vector. Its state is:

\begin{itemize}
\tightlist
\item \textbf{Experience history} \(\mathcal{H}_b^t = \{(\mathbf{x}_i, \mathbf{x}_j, q_{ij})\}\): (party 1 type, party 2 type, realized match output) triples from mediated matches (§3c). The party labels have no directional meaning because the broker's pair features are symmetric.
\item \textbf{Roster} \(\text{Roster}^t\): the broker's standing access base (§6b).
\item \textbf{Current clients} \(D^t\): agents who outsource to the broker in period \(t\) (§6b).
\item \textbf{Reputation} \(\text{rep}^t\): the average broker satisfaction of current clients (§7).
\end{itemize}

\subsection{2. The Matching Problem}\label{the-matching-problem}

Match output is symmetric and pair-level. The deterministic signal has three parts: agent general quality, pair interaction, and a pair-dependent multiplier applied to the interaction.

\subsubsection{2a. Match Quality}\label{a.-match-quality}

Let \(q_{ij}\) be the per-period value of the match between agents \(i\) and \(j\), measured in monetary units:

\[q_{ij} = Q + f(\mathbf{x}_i, \mathbf{x}_j) + \varepsilon_{ij}, \qquad
\varepsilon_{ij} \sim \mathcal{N}(0, \sigma_\varepsilon^2)\]

where \(Q = 1.0\) is a constant offset, \(\sigma_\varepsilon = 0.10\), and \(\varepsilon_{ij}\) is idiosyncratic match variation.

The deterministic signal \(f: \mathbb{R}^d \times \mathbb{R}^d \to \mathbb{R}\) is unknown to all agents and fixed for the duration of the simulation:

\[f(\mathbf{x}_i, \mathbf{x}_j) = \rho \cdot \frac{1}{2}\!\left[\mathbf{x}_i^\top \mathbf{c} + \mathbf{x}_j^\top \mathbf{c}\right] + (1 - \rho) \cdot g(\mathbf{x}_i, \mathbf{x}_j) \cdot \mathbf{x}_i^\top \mathbf{A} \mathbf{x}_j\]

where \(\rho \in [0,1]\) mixes general quality (§2b) and interaction (§2c).

\subsubsection{2b. Agent General Quality}\label{b.-agent-general-quality}

\(\mathbf{c}\) is an ideal type vector. Each agent's general-quality score is \(\mathbf{x}_i^\top\mathbf{c}\).

A pair contributes the average score \(\frac{1}{2}(\mathbf{x}_i^\top\mathbf{c}+\mathbf{x}_j^\top\mathbf{c})\) to \(f\). The ideal type \(\mathbf{c}\in\mathbb{R}^d\) is drawn at initialization as a perturbation of a random point on the agent type curve:

\[u_c\sim\operatorname{Unif}[0,1],\qquad
\boldsymbol{\epsilon}_c\sim\mathcal N\!\left(\mathbf 0,\frac{\sigma_x^2}{d}\mathbf I_d\right),\qquad
\mathbf c=\mathbf x(u_c)+\boldsymbol{\epsilon}_c.\]

Unlike agent types, \(\mathbf c\) is not re-normalized after the perturbation.

\subsubsection{2c. Match-Specific Interaction}\label{c.-match-specific-interaction}

The interaction term is \(g(\mathbf{x}_i,\mathbf{x}_j)\,\mathbf{x}_i^\top\mathbf{A}\mathbf{x}_j\), where \(g\) is a pair-dependent multiplier.

\(\mathbf{A}\) is a symmetric positive definite interaction matrix. Matrix \(\mathbf{A}\) gives a smooth bilinear interaction geometry, while \(\mathbf{B}\) divides pairs into two regions of the same type space. \(\mathbf{A}\) and \(\mathbf{B}\) are fixed for the duration of the simulation.

\textbf{Base interaction.} Draw \(\mathbf{M}_A\in\mathbb{R}^{d\times d}\) with iid \(\mathcal{N}(0,1)\) entries and set

\[
\mathbf{A}=\mathbf{M}_A^\top\mathbf{M}_A\cdot\frac{d}{\operatorname{tr}(\mathbf{M}_A^\top\mathbf{M}_A)}.
\]

Thus \(\mathbf{A}\) is symmetric positive definite and \(\operatorname{tr}(\mathbf{A})=d\).

For an isotropic unit vector \(\mathbf{x}\), \(E[\mathbf{x}\mathbf{x}^\top]=\mathbf{I}/d\), so \(E[\mathbf{x}^\top\mathbf{A}\mathbf{x}]=\operatorname{tr}(\mathbf{A})/d=1\). This normalization fixes the average interaction scale.

\textbf{Difficulty parameter.} A second symmetric matrix \(\mathbf{B}\in\mathbb{R}^{d\times d}\) defines

\[g(\mathbf{x}_i, \mathbf{x}_j) = 1 + \delta \cdot \text{sign}(\mathbf{x}_i^\top \mathbf{B} \mathbf{x}_j)\]

where \(\delta\in[0,1]\) (default 0.5).

The multiplier depends only on the sign of \(\mathbf{x}_i^\top\mathbf{B}\mathbf{x}_j\). Positive gives \(1+\delta\), negative gives \(1-\delta\), and zero gives 1. At \(\delta=0\), the boundary between the two regions has no effect. As \(\delta\) increases, crossing it changes the interaction more sharply, which makes the complementarity component harder to learn. Because \(\mathbf{B}\) is symmetric, the multiplier is symmetric in \(i\) and \(j\).

The regime matrix \(\mathbf{B}\) is constructed so the regime split is not a rescaling of the base interaction on the realized type distribution. Let

\[\boldsymbol{\Sigma}_x = \frac{1}{N} \sum_{i=1}^N \mathbf{x}_i \mathbf{x}_i^\top\]

be the uncentered second-moment matrix of the realized type vectors.

To draw the raw regime matrix, draw \(\mathbf{Z}\in\mathbb{R}^{d\times d}\) with iid \(\mathcal{N}(0,1)\) entries, symmetrize it as \(\mathbf{H}=(\mathbf{Z}+\mathbf{Z}^\top)/2\), and remove its average diagonal value:

\[
\mathbf{B}_0 = \mathbf{H} - \frac{\operatorname{tr}(\mathbf{H})}{d}\mathbf{I}_d.
\]

This makes \(\mathbf{B}_0\) symmetric with zero trace.

Under the weighted inner product defined as

\[\langle \mathbf{M}, \mathbf{N} \rangle_{\boldsymbol{\Sigma}_x} = \operatorname{tr}(\boldsymbol{\Sigma}_x \mathbf{M} \boldsymbol{\Sigma}_x \mathbf{N}),\]

remove the projection of \(\mathbf{B}_0\) onto \(\mathbf{A}\) and normalize:

\[\mathbf{B}_{\text{raw}} = \mathbf{B}_0 - \frac{\operatorname{tr}(\boldsymbol{\Sigma}_x \mathbf{B}_0 \boldsymbol{\Sigma}_x \mathbf{A})}{\operatorname{tr}(\boldsymbol{\Sigma}_x \mathbf{A} \boldsymbol{\Sigma}_x \mathbf{A})} \mathbf{A}, \qquad \mathbf{B} = \frac{\mathbf{B}_{\text{raw}}}{\lVert \mathbf{B}_{\text{raw}} \rVert_F}.\]

This gives \(\operatorname{tr}(\boldsymbol{\Sigma}_x\mathbf{B}\boldsymbol{\Sigma}_x\mathbf{A})=0\) and \(\|\mathbf{B}\|_F=1\).

The weighted projection makes \(\mathbf{B}\) orthogonal to \(\mathbf{A}\) with respect to the realized type distribution, rather than with respect to unused directions in \(\mathbb{R}^d\). In this context, orthogonality means

\[E\!\left[(\mathbf{x}_i^\top\mathbf{A}\mathbf{x}_j)(\mathbf{x}_i^\top\mathbf{B}\mathbf{x}_j)\right] = \operatorname{tr}(\boldsymbol{\Sigma}_x\mathbf{A}\boldsymbol{\Sigma}_x\mathbf{B}) = 0,\]

where the expectation is over independent draws from the realized type distribution. Because agent types lie near a curve, this removes overlap between the two bilinear scores where agents actually occur.

\(\mathbf{B}\) remains symmetric after projection and normalization.

Symmetry of \(\mathbf{A}\) and \(\mathbf{B}\) makes \(f(\mathbf{x}_i,\mathbf{x}_j)=f(\mathbf{x}_j,\mathbf{x}_i)\).

\subsection{3. Learning}\label{learning}

Before searching for matches, agents and the broker may update their prediction models on their accumulated histories. Updated models are used to score candidates (§6). The broker is eligible to refit every period. Agent \(i\) is eligible when \(i\bmod 2=t\bmod 2\), so each agent is eligible every other period. An eligible model refits only when its history is nonempty and it has received new observations since its previous fitting call. A parameter switch selects the prediction model. Neural-network learning is the default. Ridge regression is the alternate experiment.

\subsubsection{3a. Prediction Models and Fitting}\label{a.-architecture-and-fitting}

In the default model, both agents and the broker use a fully-connected network with one hidden layer, ReLU activations, and a single linear output:

\[\hat{q}(\mathbf{z}) = \mathbf{w}_2^\top \text{ReLU}(\mathbf{W}_1 \mathbf{z} + \mathbf{b}_1) + b_2\]

where \(\mathbf{z}\) is the input feature vector and the remaining symbols are network weights and biases. Agents receive raw partner type vectors, the broker receives pairs of them.

\textbf{Training rule.} Network weights minimize MSE with the Adam optimizer (learning rate \(\eta_{\mathrm{lr}}\), default 0.01) on the current training window. No explicit weight regularization is applied. At \(t = 0\), each network sets \(\mathbf W_1\) using He initialization, \(\mathbf b_1=\mathbf 0\), \(\mathbf w_2=\mathbf 0\), and \(b_2=Q\) (§2a). An untrained predictor therefore returns exactly \(Q\) for every input. Networks with seed observations then take \(E_{\text{init}}\) gradient steps (default 200) on their seed histories.

\textbf{Alternate Ridge rule.} In the Ridge experiment, agents and the broker instead use

\[
\hat q_i(\mathbf x_j)=\alpha_i+\boldsymbol\beta_i^\top\mathbf x_j,
\qquad
\hat q_b(\mathbf x_i,\mathbf x_j)=\alpha_b+\boldsymbol\gamma_b^\top\psi(\mathbf x_i,\mathbf x_j).
\]

Agent and broker fits minimize mean squared error plus \(\lambda_a\) and \(\lambda_b\), respectively, times the squared Euclidean norm of the slope coefficients. Intercepts are not penalized. Raw type and pair features are used without sample standardization. Before its first fit, a Ridge predictor has zero slopes and intercept \(Q\). An initial joint baseline calibration applied the same seven candidate penalties from \(10^{-4}\) to 1 to both learners over ten fixed seeds separate from the reporting seeds. Maximizing the median across seeds of the period 401--500 mean of the two holdout rank correlations selected a common penalty of 0.001, which we retain as \(\lambda_b\). Holding \(\lambda_b\) fixed, a separate agent calibration compared \(\lambda_a\in\{0.001,0.003,0.01,0.03,0.1,0.3,0.5\}\) on the same seeds and selected \(\lambda_a=0.003\) by the median late-period agent holdout rank correlation.

Subsequent NN updates warm-start from the previous weights and persistent Adam moment state. The number of gradient steps adapts to the ratio of new observations to total history:

\[E_t = \max\!\left(E_{\text{steps}}, \; \left\lceil E_{\text{init}} \cdot \frac{n_{\text{new}}}{n_{\text{total}}} \right\rceil\right)\]

where \(n_{\text{new}}\) is the number of observations added since that model's previous fitting call, \(n_{\text{total}} = |\mathcal{H}^t|\) is full history size, and \(E_{\text{steps}}\) (default 100) is the minimum per-call step count. For an agent, \(n_{\text{new}}\) can span multiple periods because agents fit on an alternating schedule. Agents and the broker use the same NN step-count rule. Ridge is solved directly at each fitting call.

Initialization is period 0 for each learner. Its seed-history count is stored as the first period boundary. Each later boundary stores cumulative history at the end of a simulation period. A fitting call uses observations from the most recent \(W_p = 40\) completed simulation periods. Before 40 simulation periods are available, the window also includes period 0. Once 40 are available, the period-0 seed history is excluded. If the selected history block exceeds \(M_{\text{train}} = 2000\) observations, fitting deterministically keeps 2000 spaced observations across that block, including its oldest and newest observations. If the cap is one, it keeps the newest observation. The cap changes only the fitting sample, not the stored history or the period window. NN and Ridge use the same window, cap, fitting cadence, and histories.

\subsubsection{\texorpdfstring{3b. Agent \(i\)'s Model}{3b. Agent i's Model}}\label{b.-agent-is-model}

\textbf{History.} \(\mathcal{H}_i^t = \{(\mathbf{x}_j, q_{ij})\}_{m=1}^{n_i}\) records the other party's type and realized match output from every match involving \(i\).

\textbf{Input.} The agent's network takes the partner's type as input: \(\mathbf{z} = \mathbf{x}_j\).

The hidden width is \(h_A = 2d\), so the baseline uses 16 hidden units.

\subsubsection{3c. Broker's Model}\label{c.-brokers-model}

\textbf{History.} \(\mathcal{H}_b^t = \{(\mathbf{x}_i, \mathbf{x}_j, q_{ij})\}_{m=1}^{n_b}\) records both parties' types and realized match output from broker-mediated matches.

\textbf{Input.} The broker's network takes a symmetric pair feature vector:

\[
\psi(\mathbf{x}_i,\mathbf{x}_j)
= \left[
\mathbf{x}_i+\mathbf{x}_j;\;
\operatorname{vech}\!\left(\frac{\mathbf{x}_i\mathbf{x}_j^\top+\mathbf{x}_j\mathbf{x}_i^\top}{2}\right)
\right].
\]

The half-vectorization \(\operatorname{vech}\) keeps the lower-triangular entries of the symmetric matrix, so the broker input dimension is \(d_B = d + d(d+1)/2\), equal to 44 at \(d = 8\).

The broker hidden width is \(h_B = 8d\), so the baseline uses 64 hidden units.

Because \(\psi(\mathbf{x}_i,\mathbf{x}_j)=\psi(\mathbf{x}_j,\mathbf{x}_i)\), each observed unordered pair contributes one training row.

The broker learns only from mediated matches. It does not observe self-search outcomes.

\subsubsection{3d. Ridge Broker Ablations}\label{d.-ridge-broker-ablations}

Three separate Ridge-only ablations change the broker while leaving every agent's local Ridge model unchanged. \textbf{Size-matched pooled history} uniformly samples the broker's windowed and capped fitting data without replacement, using the smaller of its available sample size and the nearest-integer median sample size among agents with nonempty histories, with half-integers rounded up. \textbf{Single-principal history} randomly retains one endpoint when each broker observation is recorded, fits \(g_b(\mathbf x)=\alpha_b+\boldsymbol\gamma_b^\top\mathbf x\), and scores a pair as \(g_b(\mathbf x_i)+g_b(\mathbf x_j)-\bar q_b\), where \(\bar q_b\) is the current fitting-sample output mean and defaults to \(Q\) before fitting. \textbf{Additive Ridge} uses \(\mathbf x_i+\mathbf x_j\) without bilinear features. All use the calibrated agent and broker penalties, period window, observation cap, and fitting cadence.

\subsection{4. Match Economics}\label{match-economics}

Matches use a common reservation threshold and channel-specific cost accounting.

At initialization, 10,000 random agent pairs define the calibration mean \(\bar{q}_{\text{cal}} = E[q]\). This calibration quantity scales \(r\), \(\phi\), and \(c_s\).

\subsubsection{4a. Outside Options}\label{a.-outside-options}

All agents share a reservation threshold \(r\), calibrated at initialization:

\[r = f_r \cdot \bar{q}_{\text{cal}}\]

where \(\bar{q}_{\text{cal}}\) is the calibration mean and \(f_r\) is the reservation coefficient (default 0.60). Since \(\phi\) and \(c_s\) do not depend on \(r\) (§4c), \(f_r\) may take any nonnegative value, including \(f_r \geq 1\), and is swept as an analysis axis.

\subsubsection{4b. Participation Constraints}\label{b.-participation-constraints}

Demanders emit offers only above the reservation threshold. Self-search uses the demander evaluation \(\hat{q}_{i}(\mathbf{x}_j)>r\). Brokered search uses the broker pair prediction \(\hat{q}_b(\{i,j\})>r\). A one-sided offer is accepted only if the receiver's evaluation exceeds \(r\); reciprocal offers are accepted automatically (§6d).

\subsubsection{4c. Search Costs}\label{c.-search-costs}

The self-search cost and broker fee share one calibrated level:

\[
\phi = c_s = \lambda_c \cdot \bar{q}_{\text{cal}}.
\]

\(\lambda_c\) is the shared cost rate (default 0.05). The two symbols differ in payoff accounting:

\begin{itemize}
\tightlist
\item \textbf{Self-search} costs \(c_s\) on each requested self-search position, whether or not it is filled.
\item \textbf{Brokerage} charges \(\phi\) only on successful brokered placements.
\end{itemize}

\subsection{5. Network Structure and Turnover}\label{network-structure-and-turnover}

Agents interact through a single undirected network \(G\) that determines search opportunities and structural position.

\subsubsection{5a. Network Initialization}\label{a.-network-initialization}

\(G\) is initialized as a Watts-Strogatz small-world graph on the \(N\) agents (Watts \& Strogatz, 1998), with ring degree \(k_G = 6\), rewiring probability \(p_{\text{rewire}} = 0.1\), and random agent order on the ring.

The broker is added as a permanent node with no type vector. Broker-node edges are synchronized to standing roster members, current-period broker clients, and agents currently engaged in broker-channel matches (§6b).

\subsubsection{5b. Match Tie Formation}\label{b.-match-tie-formation}

Each realized match adds an undirected \(i\)-\(j\) edge if one does not already exist. Ties persist until either node exits.

\subsubsection{5c. Agent Turnover}\label{c.-agent-turnover}

Agents exit independently each period with probability \(\eta_{\mathrm{exit}}\) (default 0.02). An exiting agent is removed everywhere. Its node slot is reused by an entrant with a fresh type drawn from the same type process (§1). The entrant receives \(\lfloor k_G/2 \rfloor\) edges to existing agents sampled with probability proportional to \(\exp(-\|\mathbf{x}_{i'} - \mathbf{x}_j\|^2)\).

\subsection{6. Search}\label{search}

At the start of each period, each agent draws active relationship demand

\[d_i \sim \text{Binomial}(K,\; p_{\text{demand}}),\]

where \(K\) is the maximum number of active demands the agent can initiate in the period and \(p_{\text{demand}}\) defaults to 0.50. If \(d_i > 0\), the agent chooses \textbf{one channel for the batch} (§7): self-search or broker.

\subsubsection{6a. Self-Search}\label{a.-self-search}

Agent \(i\)'s self-search candidate pool has two components:

\textbf{Neighbors.} For a direct network neighbor \(j\) with prior pair-specific experience, agent \(i\) uses \(\bar{q}_{ij} = \frac{1}{n_{ij}} \sum q_{ij}^{(m)}\), where \(n_{ij}\) is the number of times \(i\) and \(j\) have matched. If the graph edge exists but \(n_{ij}=0\), the agent instead uses \(\hat q_i(\mathbf x_j)\). Initial histories seed one observation for every initial graph edge. Unobserved neighbor edges arise after turnover because entrants attach structurally before matching.

\textbf{Strangers.} Each self-searching demander \(i\) independently samples a pool \(U_i^t\) of \(\min(n_{\mathrm{strangers}}, N)\) agents uniformly without replacement, where \(n_{\mathrm{strangers}} = 10\) (default). The demander evaluates pool members that are not current neighbors using \(\hat{q}_i(\mathbf{x}_j)\) (§3b).

Agent \(i\) orders feasible self-search candidates by these evaluations and emits up to \(d_i\) directed offers to candidates above \(r\). Equal evaluations are ordered uniformly at random, so an agent whose predictions are constant produces a random ranking rather than an undefined ordering.

\subsubsection{6b. Broker Roster and Access}\label{b.-broker-roster-and-access}

The broker maintains a \textbf{roster} of agents it can propose as counterparties when mediating matches. In each period, it also observes the current-client set \(D^t\) of agents who outsourced. Broker access for the period is:

\[\mathcal{A}_b^t = \text{Roster}^t \cup D^t.\]

Current clients expand broker access for the period but do not become standing roster members.

\textbf{Initialization.} The roster is seeded with a target size

\[R^* = \lceil \alpha_R N \rceil,\]

where the broker roster share \(\alpha_R\) defaults to 0.20. The model draws \(R^*\) agents uniformly at random (default 200 at \(N = 1000\)). The broker's history is seeded from existing roster-roster ties in the initial graph, capped at 100 observed ties.

\textbf{Roster refresh.} At the start of each period, each roster member independently exits with probability \(p_{\text{roster}}\) (default \(0.02\)). If \(\widetilde{\text{Roster}}^t\) is the post-churn roster,

\[\widetilde{\text{Roster}}^t = \{i \in \text{Roster}^{t-1} : u_i^t > p_{\text{roster}}\}, \qquad u_i^t \overset{iid}{\sim} \mathrm{Unif}[0,1],\]

then the broker samples without replacement from \(\{1,\ldots,N\}\setminus \widetilde{\text{Roster}}^t\) until \(|\text{Roster}^t| = R^*\), or until the population is exhausted. Broker-node edges are synchronized to the standing roster, current clients, and current broker-mediated matches (§5a).

\subsubsection{6c. Broker-Mediated Search}\label{c.-broker-mediated-search}

For every unordered pair \(\{i,j\}\) such that one side is a broker-client demander and the other side is in \(\mathcal{A}_b^t\), the broker computes predicted match quality. For each broker-client demander \(i\), it ranks accessible counterparties above \(r\), breaking equal predictions uniformly at random, and keeps the highest-ranked candidates up to \(i\)'s residual active broker demand. Directed broker offers are then processed in a global score ranking with a new, independent random ordering within score ties. If both sides are broker-client demanders with remaining active demand, the same unordered pair can generate reciprocal broker offers.

\subsubsection{6d. Offer Acceptance and Match Formation}\label{d.-offer-acceptance-and-match-formation}

All self-search and broker offers enter one shared market, grouped by unordered pair.

\begin{enumerate}
\def\labelenumi{\arabic{enumi}.}
\tightlist
\item If pair \(\{i,j\}\) contains reciprocal offers \(i \to j\) and \(j \to i\), the relationship is accepted automatically.
\item If pair \(\{i,j\}\) contains one directed offer \(i \to j\), receiver \(j\) evaluates \(i\) using \(\bar{q}_{ji}\) for known partners and \(\hat{q}_j(\mathbf{x}_i)\) otherwise. The offer is accepted iff this value exceeds \(r\).
\item Each accepted unordered pair forms at most one current-period relationship.
\end{enumerate}

\textbf{At market finalization, for each accepted match:} 1. Realized output is drawn: \(q_{ij} = Q + f(\mathbf{x}_i, \mathbf{x}_j) + \varepsilon_{ij}\). 2. Both parties add the observation to their histories: \(i\) adds \((\mathbf{x}_j, q_{ij})\) to \(\mathcal{H}_{i}\); \(j\) adds \((\mathbf{x}_i, q_{ij})\) to \(\mathcal{H}_{j}\). 3. If either directed offer used the broker channel, the broker adds \((\mathbf{x}_i, \mathbf{x}_j, q_{ij})\) to \(\mathcal{H}_b\) once. 4. An edge is added between \(i\) and \(j\) in \(G\) (if not already present).

A demander can have at most \(d_i\) outgoing offers, but can receive any number of acceptable incoming offers. Thus \(K\) is an active-demand bound, not a capacity constraint on receiving offers.

Before the next period begins, match lists are cleared. The next period begins with no active relationships.

\subsection{7. The Outsourcing Decision}\label{the-outsourcing-decision}

\subsubsection{7a. Satisfaction Tracking}\label{a.-satisfaction-tracking}

Each agent \(i\) maintains a satisfaction index \(s_{i,c}^t\) for each search channel \(c \in \{\text{self}, \text{broker}\}\). For the chosen channel, the index updates once per active-demand period as an EWMA of realized match value net of search costs:

\[s_{i,c}^{t+1} = (1 - \omega)\,s_{i,c}^t + \omega \cdot \tilde{q}\]

where \(\omega = 0.2\) and \(\tilde{q}\) is the period input in the table. The denominator is requested demand \(d_i\), so unfilled requests contribute zero output. A reciprocal relationship can update two channel-specific states, one for each directed offer.

{
\begin{longtable}[]{@{}
  >{\raggedright\arraybackslash}p{(\linewidth - 2\tabcolsep) * \real{0.2250}}
  >{\raggedright\arraybackslash}p{(\linewidth - 2\tabcolsep) * \real{0.7750}}@{}}
\toprule\noalign{}
\begin{minipage}[b]{\linewidth}\raggedright
Channel
\end{minipage} & \begin{minipage}[b]{\linewidth}\raggedright
Satisfaction input \(\tilde{q}\)
\end{minipage} \\
\midrule\noalign{}
\endhead
\bottomrule\noalign{}
\endlastfoot
Self-search & \(\dfrac{\sum q_{ij}}{d_i} - c_s\), summing over accepted directed offers made by \(i\) through self-search \\
Brokered & \(\dfrac{\sum q_{ij} - \phi \cdot n_{i,\text{broker success}}}{d_i}\), summing over accepted directed offers made by \(i\) through the broker \\
\end{longtable}
}

Brokered failure gives \(\tilde{q}=0\), self-search failure gives \(\tilde{q}=-c_s\).

\(s_{i,\text{self}}^t\) covers all self-search, including known neighbors.

At initialization, self-satisfaction is the mean of seed match outcomes and broker-satisfaction is the broker's seed-data reputation (§7c). Fresh entrants set self-satisfaction to mean neighbor self-satisfaction and broker-satisfaction to current broker reputation.

The first broker choice, regardless of placement success, switches the agent from broker reputation to its own \(s_{i,\text{broker}}^t\) (§7b).

\subsubsection{7b. Decision Rule}\label{b.-decision-rule}

Each agent with \(d_i > 0\) compares:

\begin{itemize}
\tightlist
\item \textbf{\(\text{score}_{\text{self}}\)} \(= s_{i,\text{self}}^t\)
\item \textbf{\(\text{score}_{\text{broker}}\)} \(= s_{i,\text{broker}}^t\) if the agent has tried the broker, otherwise \(\text{rep}_b^t\).
\end{itemize}

The agent chooses the higher-score channel. Ties are resolved by fair coin flip.

\subsubsection{7c. Broker Reputation}\label{c.-broker-reputation}

\[\text{rep}_b^{t+1} = \begin{cases} \frac{1}{|D^t|} \sum_{i \in D^t} s_{i,\text{broker}}^{t+1} & \text{if } D^t \neq \emptyset \\[4pt] \text{rep}_b^{t} & \text{otherwise} \end{cases}\]

where \(D^t\) is the set of agents who outsourced to the broker this period. With current clients, reputation is their mean post-update broker satisfaction; with no clients, it is unchanged. Reputation is initialized from the mean broker seed outcome.

\subsection{8. Period Measures}\label{period-measures}

Period measures are evaluated after current-period matching, satisfaction, and reputation updates, but before agent entry/exit turnover. Agent decisions do not depend on these measures.

\subsubsection{Network Measures}\label{network-measures}

Broker network-position measures use \(G\), including the broker node (§5a). They are computed when \(t \bmod \Delta_{\mathrm{net}}=0\) (default 20). Other period measures are recorded each period.

\textbf{Betweenness centrality.} Freeman betweenness (Freeman, 1977) for the broker node:

\[\mathrm{BC}_b = \frac{1}{(n-1)(n-2)} \sum_{u \neq b} \sum_{v \neq u} \frac{\sigma_{uv}(b)}{\sigma_{uv}}\]

where \(n = N+1\), \(\sigma_{uv}\) is the number of shortest paths from \(u\) to \(v\), and \(\sigma_{uv}(b)\) is the number passing through \(b\). The double sum counts undirected pairs from both directions, so the denominator is \((n-1)(n-2)\) (Brandes, 2001, p.~9).

\textbf{Burt's constraint.} Broker ego-network constraint (Burt, 1992):

\[\mathrm{Constraint}_b = \sum_j \left(p_{bj} + \sum_{h \neq b,j}
p_{bh}\, p_{hj}\right)^2\]

where \(p_{bj} = 1/d_b\) and \(p_{hj} = 1/d_h\) (Everett \& Borgatti, 2020). The indirect term uses intermediary \(h\)'s degree.

\textbf{Effective size.} Non-redundant broker contacts using the Borgatti (1997) binary-undirected simplification:

\[\text{ES}_b = d_b - \frac{2t_b}{d_b},\]

where \(d_b\) is the broker's degree and \(t_b\) is the number of ties among the broker's neighbors, excluding ties to the broker.

\subsubsection{Prediction Quality}\label{prediction-quality}

Prediction quality is evaluated on holdout pairs and accepted directed offers:

{
\begin{longtable}[]{@{}
  >{\raggedright\arraybackslash}p{(\linewidth - 4\tabcolsep) * \real{0.1951}}
  >{\raggedright\arraybackslash}p{(\linewidth - 4\tabcolsep) * \real{0.3415}}
  >{\raggedright\arraybackslash}p{(\linewidth - 4\tabcolsep) * \real{0.4634}}@{}}
\toprule\noalign{}
\begin{minipage}[b]{\linewidth}\raggedright
Sample
\end{minipage} & \begin{minipage}[b]{\linewidth}\raggedright
Construction
\end{minipage} & \begin{minipage}[b]{\linewidth}\raggedright
Reported measures
\end{minipage} \\
\midrule\noalign{}
\endhead
\bottomrule\noalign{}
\endlastfoot
Holdout pairs & Up to 100 agents with non-empty histories; for each sampled agent \(i\), \(\min(40,N-1)\) non-self partners \(j\). Agent and broker predictions are evaluated against noiseless \(Q + f(\mathbf{x}_i,\mathbf{x}_j)\). Both models use the same deterministic sample conditional on seed and period. & Per-agent \(R^2\), RMSE, bias, and Spearman rank correlation, averaged across sampled agents. \\
Selected offers & Accepted directed offers by channel, self-search or brokered. A reciprocal relationship can represent two directed decisions. & Channel-specific \(R^2\), RMSE, bias, and Spearman rank correlation. \\
\end{longtable}
}

The definitions are

\[
R^2 = 1 - \frac{\mathrm{MSE}}{\operatorname{Var}_{\mathrm{pop}}(q)}, \qquad
\mathrm{RMSE} = \sqrt{\frac{1}{n}\sum(\hat{q}-q)^2}, \qquad
\mathrm{Bias} = \frac{1}{n}\sum(\hat{q}-q).
\]

Rank correlation evaluates the strict ranking produced from predicted output. Candidates are ordered by predicted output, with each set of equal predictions ordered uniformly at random. Agent and broker tie-breaking draws are independent. Spearman's \(\rho_S\) is then computed between this produced ranking and the realized-output ranks; equal realized outputs retain their average ranks. The same rule is applied to each holdout candidate set and to the pooled accepted directed offers within each selected-offer channel. Thus a constant predictor produces a random ranking with expected rank correlation zero. For selected-offer summaries, \(R^2\), bias, and rank correlation are undefined when \(n < 5\) or \(\operatorname{Var}_{\mathrm{pop}}(q) < \sigma_\varepsilon^2/6\), while RMSE is defined whenever \(n>0\). Holdout summaries use common valid-agent support for all four measures: an agent-level holdout set that fails the sample-size or variance condition is excluded from the \(R^2\), RMSE, bias, and rank-correlation means.

\subsubsection{Other Measures}\label{other-measures}

{
\begin{longtable}[]{@{}
  >{\raggedright\arraybackslash}p{(\linewidth - 2\tabcolsep) * \real{0.4286}}
  >{\raggedright\arraybackslash}p{(\linewidth - 2\tabcolsep) * \real{0.5714}}@{}}
\toprule\noalign{}
\begin{minipage}[b]{\linewidth}\raggedright
Measure
\end{minipage} & \begin{minipage}[b]{\linewidth}\raggedright
Definition
\end{minipage} \\
\midrule\noalign{}
\endhead
\bottomrule\noalign{}
\endlastfoot
Access vs.~assessment & For each accepted broker-directed offer, record whether receiver \(j\) was a direct neighbor of sender \(i\) in \(G\) at the time of the offer. Non-neighbor receivers indicate access; neighbor receivers indicate assessment. \\
Match quality by channel & Average realized match output \(\bar{q}_c^t\), where \(c \in \{\text{self}, \text{brokered}\}\). \\
Mean channel satisfaction & Cross-agent means \(N^{-1}\sum_i s_{i,\text{self}}^t\) and \(N^{-1}\sum_i s_{i,\text{broker}}^t\). \\
Outsourcing rate & Outsourced requested positions divided by total requested positions. \\
Broker access size & \(\|\mathcal{A}_b^t\| = \|\text{Roster}^t \cup D^t\|\), measured after outsourcing decisions form \(D^t\) and before entry/exit. \\
Counterparty concentration & Median and maximum current-period counterparties per agent. \\
Agent-node degree summaries & Mean, median, minimum, and maximum degree, excluding the broker node. \\
\end{longtable}
}

\section{Part II. Baseline Values and Reproducibility}\label{part-ii.-baseline-values-and-reproducibility}

\subsection{9. Baseline Values and Reproducibility}\label{baseline-values-and-reproducibility}

\subsubsection{9a. Baseline Parameter Values}\label{a.-baseline-parameter-values}

The table lists baseline values only. Entries are arranged in two columns; each column follows the document order, with the run length listed first.

The baseline learning function is NN. Ridge and its broker variants are alternate experiments (§3).

{
\begin{center}
\small
\setlength{\tabcolsep}{3pt}
\begin{tabular}{@{}lll@{\hspace{1.6em}}lll@{}}
\toprule
Symbol & Baseline & Defined in & Symbol & Baseline & Defined in \\
\midrule
\(T\) & 500 & §1 & \(f_r\) & 0.60 & §4a \\
\(N\) & 1000 & §1 & \(\lambda_c\) & 0.05 & §4c \\
\(d\) & 8 & §1 & \(k_G\) & 6 & §5a \\
\(d_\gamma\) & 8 & §1 & \(p_{\text{rewire}}\) & 0.1 & §5a \\
\(\sigma_x\) & 0.5 & §1 & \(\eta_{\mathrm{exit}}\) & 0.02 & §5c \\
\(Q\) & 1.0 & §2a & \(K\) & 5 & §6 \\
\(\rho\) & 0.50 & §2a & \(p_{\text{demand}}\) & 0.50 & §6 \\
\(\delta\) & 0.5 & §2c & \(n_{\mathrm{strangers}}\) & 10 & §6a \\
\(\sigma_\varepsilon\) & 0.10 & §2a & \(\alpha_R\) & 0.20 & §6b \\
\(\eta_{\mathrm{lr}}\) & 0.01 & §3a & \(p_{\text{roster}}\) & 0.02 & §6b \\
\(E_{\text{init}}\) & 200 & §3a & \(\omega\) & 0.2 & §7a \\
\(E_{\text{steps}}\) & 100 & §3a & \(\Delta_{\mathrm{net}}\) & 20 & §8 \\
\(W_p\) & 40 periods & §3a & \(M_{\text{train}}\) & 2000 & §3a \\
\bottomrule
\end{tabular}
\end{center}
}

\subsubsection{9b. Reproducibility}\label{b.-reproducibility}

All model-event randomness flows from a single integer seed. The seed determines type draws, \(G\), matching function parameters (\(\mathbf{c}\), \(\mathbf{A}\), \(\mathbf{B}\)), broker seed roster, standing-roster churn and replenishment, random ordering within market score ties, and subsequent model events. Holdout samples and their independent agent and broker tie-breaking draws are deterministic conditional on seed and period. Simulations are fully reproducible given parameters and seed.

\section{References}\label{references}

Brandes, U. (2001). A faster algorithm for betweenness centrality. \emph{Journal of Mathematical Sociology}, \emph{25}(2), 163--177.

Borgatti, S. P. (1997). Structural holes: Unpacking Burt's redundancy measures. \emph{Connections}, \emph{20}(1), 35--38.

Burt, R. S. (1992). \emph{Structural holes: The social structure of competition}. Harvard University Press.

Everett, M. G., \& Borgatti, S. P. (2020). Unpacking Burt's constraint measure. \emph{Social Networks}, \emph{62}, 50--57.

Freeman, L. C. (1977). A set of measures of centrality based on betweenness. \emph{Sociometry}, \emph{40}(1), 35--41.

Watts, D. J., \& Strogatz, S. H. (1998). Collective dynamics of `small-world' networks. \emph{Nature}, \emph{393}(6684), 440--442.

\end{document}
